\chapter{Introduction}
\label{ch:intro}

\subsubsection*{New Physics without New Particles}

The discovery of the Higgs boson at the Large Hadron Collider (LHC)~\cite{ATLAS:2012yve,CMS:2012qbp}
completed the particle content of the Standard Model (SM), and the measurements that followed have
confirmed its predictions in every sector tested so far. Moreover, no elementary particle beyond
those of the SM has been observed, either at the LHC or in the broad range of
lower-energy experiments that probe the same interactions indirectly. However, this success coexists
with a list of questions that the SM leaves unanswered. Some are observational, such as
the nature of dark matter~\cite{Bertone:2004pz,Cirelli:2024ssz}, the origin of neutrino
masses~\cite{Super-Kamiokande:1998kpq,SNO:2002tuh}, and the asymmetry between matter and
antimatter~\cite{Sakharov:1967dj,Planck:2018vyg}. Others concern the structure of the theory itself,
such as the smallness of the CP-violating angle of quantum chromodynamics (QCD), the hierarchical
pattern of fermion masses and mixings, and the stability of the electroweak scale against large
quantum corrections~\cite{tHooft:1979rat}. Therefore, it is widely held that the SM is
the low-energy description of a more fundamental theory, whose new degrees of freedom are either too
heavy to be produced at the energies explored so far, at a scale $\Lambda$, or too weakly coupled to
have been observed.

If that scale lies beyond direct reach, the new dynamics can only be seen through its indirect
effects, and these take a universal form. Indeed, at energies $E \ll \Lambda$ the heavy states
cannot be produced, and their exchange reduces to local interactions among the light fields,
suppressed by powers of $E/\Lambda$. Whatever the underlying theory is, its low-energy imprint is
therefore captured, at any given order in $E/\Lambda$, by a finite number of coefficients
multiplying higher-dimensional operators built out of the light fields. The search for physics
beyond the SM is thereby recast as the determination of these coefficients, and the
effective field theory (EFT) that organizes them becomes the common language in which collider
measurements, low-energy precision tests, and model building can be compared.

This language is useful only if two requirements are met. On the one hand, the coefficients are
generated at the scale $\Lambda$ and measured at much lower energies, often several orders of
magnitude apart, so their evolution between the two scales must be under control. On the other hand,
the coefficients cannot take arbitrary values, because the EFT must remain consistent at the
energies where it is applied and must admit a UV completion with a unitary, causal, and local
$S$-matrix, which restricts the region they may occupy before any measurement is made. This thesis
is concerned with both requirements. Its second part computes how the coefficients evolve, through
the anomalous dimensions of the corresponding operators, and studies the structure of operator mixings. Its third part determines which regions of coefficient space are compatible
with the unitarity, analyticity, and locality of the $S$-matrix. Both questions are addressed mainly
for the two EFTs most directly connected to the open problems listed above, the Standard Model
Effective Field Theory (SMEFT) and the EFT of axion-like particles (ALPs), and both are approached
with the same set of tools, the on-shell methods for scattering amplitudes.

\subsubsection*{The Effective-Field-Theory Paradigm}

The logic of EFTs is older than the SM, as the Fermi and Euler--Heisenberg theories
show~\cite{Fermi:1934hr,Heisenberg:1936nmg}, and is not specific to
it~\cite{Weinberg:1978kz,Georgi:1993mps,Manohar:2018aog}. It rests on the separation of scales,
which is ultimately what makes physics possible. Nature spans an enormous range of distances and
energies, and yet each phenomenon can be described without knowing what happens at much shorter
distances. Indeed, the motion of the planets can be computed without knowing their internal
composition, the spectrum of the hydrogen atom without knowing the structure of the proton, and
chemistry without knowing QCD. In each case, the short-distance dynamics enters through a handful of
parameters, such as masses and charges, while its remaining effects are suppressed by powers of the
ratio of the two scales. For instance, the finite size of the proton shifts the $S$ levels of
hydrogen, relative to their binding energy, only at order $(r_p/a_0)^2\sim10^{-10}$, where $r_p$ is
the proton radius and $a_0$ the Bohr radius. If this were not the case, no prediction could be made
without first knowing the ultimate theory of nature. 

In quantum field theory, the same property
takes the form of the decoupling theorem, according to which, when the heavy masses can be taken
large independently of the couplings, the effects of heavy states on low-energy observables are
either absorbed into the renormalization of the couplings of the light fields or suppressed by
inverse powers of their masses~\cite{Appelquist:1974tg}. An EFT makes this statement systematic. It
retains only the degrees of freedom that are relevant at the scale of interest and encodes the effects
of the heavy ones in local interactions, organized as an expansion in the ratio of the two scales.

Concretely, an EFT is specified by three ingredients:
\begin{itemize}
    \item The degrees of freedom, namely the fields that are light at the energies of interest;
    \item The symmetries that the interactions of these fields must respect, whether exact,
    approximate, or spontaneously broken and nonlinearly realized;
    \item A set of expansion parameters, given by ratios of the low-energy scales to the heavy ones,
    such as $E/\Lambda$, together with a power counting that assigns to each operator the order at
    which it contributes.
\end{itemize}
Once these ingredients are fixed, the EFT Lagrangian contains all the local operators built out of
the light fields and compatible with the symmetries, ordered by the power counting, and, at any
finite order, only a finite number of them contribute.

This construction acquires a precise meaning in the Wilsonian approach to
renormalization~\cite{Wilson:1971bg,Wilson:1973jj,Polchinski:1983gv}. There, a quantum field theory
is always defined with a cutoff, and lowering the cutoff amounts to integrating out, in the path
integral, the modes with momenta between the old and the new cutoff. The resulting Wilsonian action
contains all the local interactions allowed by the symmetries, with couplings that flow as the
cutoff is lowered. Operators of dimension larger than four are irrelevant in this flow, since their
effects at low energies are suppressed by powers of the ratio between the energy and the scale at
which they are generated, which explains why the low-energy physics depends on the short-distance
dynamics only through a few parameters. When the heavy degrees of freedom are particles of mass $M$,
described by a field $\Phi$ coupled to the light fields $\phi$, the same procedure defines the EFT
through
\begin{equation}\label{eq:intro_pathintegral}
    e^{iS_{\rm EFT}[\phi]} = \int\mathcal D\Phi\;e^{iS[\phi,\Phi]}\,.
\end{equation}
At tree level, the path
integral over $\Phi$ is dominated by the solution $\Phi_c[\phi]$ of its equation of motion, $\delta
S/\delta\Phi|_{\Phi=\Phi_c}=0$, which is non-local on scales of order $1/M$. Expanding it in powers
of $\partial/M$ yields a series of local operators suppressed by powers of $1/M$, so that
\begin{equation}\label{eq:intro_treeloop}
    S_{\rm EFT}[\phi] = S[\phi,\Phi_c[\phi]] + \frac{i}{2}\Tr \log\!\left(-\frac{\delta^2S}{\delta\Phi\,\delta\Phi}\right)\!\Bigg|_{\Phi=\Phi_c}+\dots\,,
\end{equation}
where the second term, which arises from the Gaussian integration over the fluctuations of $\Phi$
around $\Phi_c$, gives the one-loop contributions~\cite{Henning:2014wua}, and the ellipsis denotes
higher orders. In practice, EFTs are defined in dimensional regularization with mass-independent
schemes, where no explicit cutoff appears. The heavy fields are then removed at the scale of their
mass by the matching procedure discussed below, which implements the decoupling explicitly, while
the evolution of the couplings between thresholds is described by the renormalization group of the
EFT.

The Fermi theory of the weak interaction~\cite{Fermi:1934hr}, in its later $V-A$
form~\cite{Feynman:1958ty,Sudarshan:1958vf}, is the classic example. At momentum transfers much
smaller than the mass $m_W$ of the $W$ boson, the propagator of the $W$ boson exchanged between two
fermion currents cannot be resolved and shrinks to a point:
\begin{equation}
    \begin{tikzpicture}[baseline={(0,-0.1)}]
        \begin{feynman}
            \vertex (a) at (0,0.75);
            \vertex (b) at (0,-0.75);
            \vertex (i1) at (-1.4,1.3) {\(\nu_\mu\)};
            \vertex (f1) at (1.4,1.3) {\(\mu^-\)};
            \vertex (i2) at (-1.4,-1.3) {\(e^-\)};
            \vertex (f2) at (1.4,-1.3) {\(\nu_e\)};
            \diagram*{
                (i1) -- [fermion] (a) -- [fermion] (f1),
                (i2) -- [fermion] (b) -- [fermion] (f2),
                (a) -- [boson, edge label=\(W\)] (b),
            };
        \end{feynman}
        \node at (0,-2.1) {$\displaystyle\propto\frac{g^2}{q^2-m_W^2}$};
    \end{tikzpicture}
    \qquad\xrightarrow{\quad |q^2|\,\ll\, m_W^2\quad}\qquad
    \begin{tikzpicture}[baseline={(0,-0.1)}]
        \begin{feynman}
            \vertex [dot] (c) at (0,0) {};
            \vertex (i1) at (-1.4,1.3) {\(\nu_\mu\)};
            \vertex (f1) at (1.4,1.3) {\(\mu^-\)};
            \vertex (i2) at (-1.4,-1.3) {\(e^-\)};
            \vertex (f2) at (1.4,-1.3) {\(\nu_e\)};
            \diagram*{
                (i1) -- [fermion] (c) -- [fermion] (f1),
                (i2) -- [fermion] (c) -- [fermion] (f2),
            };
        \end{feynman}
        \node at (0,-2.1) {$\displaystyle\propto-\frac{g^2}{m_W^2}$};
    \end{tikzpicture}
\end{equation}
where $q$ is the momentum of the $W$ boson and $g$ is the $SU(2)_L$ gauge coupling.
Expanding the propagator in powers of $q^2/m_W^2$, the leading term yields the four-fermion
interaction
\begin{equation}\label{eq:intro_Fermi}
    \mathcal L_{\rm Fermi} = -\frac{4G_F}{\sqrt 2}\,(\overline\mu\gamma^\alpha P_L\nu_\mu)(\overline\nu_e\gamma_\alpha P_L e)
    + \text{H.c.}\,,\qquad \frac{G_F}{\sqrt 2}=\frac{g^2}{8m_W^2}\,,
\end{equation}
where $P_L$ is the left-handed projector, while the subleading terms, suppressed by powers of
$q^2/m_W^2$, correspond to operators of higher dimension with additional derivatives. At low
energies, the underlying gauge theory is therefore visible only through the Fermi constant $G_F$ and
the chiral structure of the currents. Moreover, the Fermi theory signals its own limitations. Since
$G_F$ has the dimension of an inverse mass squared, the amplitude of $\nu_\mu e^-\to\mu^-\nu_e$
grows as $G_F s$ and violates perturbative unitarity at energies of several hundred GeV, indicating
that the theory must be modified below that scale, as it eventually was by the $W$ boson.

EFTs are used in two complementary ways. In the top-down approach, the theory at high energies is
known, and the EFT is derived from it by integrating out the heavy degrees of freedom.
This may seem superfluous, since the full theory contains all the information, but it is often the
most efficient, and sometimes the only, way to compute low-energy observables. First, a calculation
that involves widely separated scales is replaced by a sequence of simpler ones, each involving a
single scale. Second, logarithms of the ratio of the scales, which appear at every order in
perturbation theory multiplied by powers of the couplings and invalidate the fixed-order expansion
when their product is of order one, are resummed by the renormalization group of the EFT. The weak
decays of hadrons are a standard example, since the QCD corrections enhanced by logarithms of the
ratio between $m_W$ and the hadronic scale, such as $m_W/m_b$ in $B$-meson decays, are resummed
within the effective weak Hamiltonian~\cite{Gaillard:1974nj,Altarelli:1974exa,Buchalla:1995vs}.
Third, when the full theory is strongly coupled at low energies, the EFT cannot be derived in
perturbation theory, and is instead constructed from the symmetries of the full theory. This is the
case of chiral perturbation theory~\cite{Weinberg:1978kz,Gasser:1983yg,Gasser:1984gg}, which
describes the pions, kaons, and eta meson as the pseudo-Nambu--Goldstone bosons of the spontaneously
broken chiral symmetry of QCD, in an expansion in powers of momenta and pseudo-Goldstone masses over
the chiral-symmetry-breaking scale $\Lambda_\chi\sim4\pi f_\pi\sim1$~GeV~\cite{Manohar:1983md}.
Indeed, although QCD is known, its low-energy dynamics cannot be computed in perturbation theory,
and chiral perturbation theory makes predictions from symmetry alone, with the dependence on the
underlying dynamics encoded in a finite number of low-energy constants at each order, which can be
extracted from data or from lattice QCD.

In the bottom-up approach, instead, the theory at high energies is unknown, and the EFT is
constructed from the light degrees of freedom and the symmetries observed at low energies. Its
rationale is a conjecture --- often called Weinberg's folk theorem --- stated as follows \cite{Weinberg:1978kz}:
\begin{quotation}
``If one writes down the most general
possible Lagrangian, including \emph{all} terms consistent with assumed symmetry principles,
and then calculates matrix elements with this Lagrangian to any given order of perturbation theory, the result
will simply be the most general possible $S$-matrix consistent with analyticity, perturbative unitarity, cluster decomposition and the assumed symmetry principles.''
\end{quotation}
Applied to physics beyond the SM, this reasoning leads to the Lagrangian
\begin{equation}\label{eq:intro_Leff}
    \mathcal L_{\rm EFT} = \mathcal L_{\rm SM} + \sum_i \frac{c_i}{\Lambda^{[\mathcal O_i]-4}}\,\mathcal O_i\,,
\end{equation}
where the operators $\mathcal O_i$ are built out of the SM fields, have mass dimension
$[\mathcal O_i]>4$, and respect the assumed symmetries, while the dimensionless Wilson coefficients
$c_i$ encode the unknown dynamics at the scale $\Lambda$. The expansion is organized by dimensional
analysis. Indeed, an operator of dimension $d$ modifies an amplitude at the energy $E$ by a relative
amount of order $(E/\Lambda)^{d-4}$, where $E$ may be replaced by the electroweak scale or by other
low-energy scales of the process, so that to any given accuracy only a finite number of operators is
needed, and the theory, although non-renormalizable in the traditional sense, is predictive order by
order in $E/\Lambda$.

% The counting has a subtlety that matters for what follows. In most observables, the leading effect of
% a dimension-six operator is its interference with the SM amplitude, of order $c_i
% E^2/\Lambda^2$, while the square of the new contribution is of order $c_i^2E^4/\Lambda^4$, the same
% order as the interference of dimension-eight operators and of amplitudes with two insertions of
% dimension-six operators. Moreover, the interference can be suppressed, for instance by helicity
% selection rules~\cite{Azatov:2016sqh}, in which case the squared terms dominate. Whether such terms
% are retained, which distinguishes the quadratic from the linear fits of global
% analyses~\cite{Brivio:2019ius,Ellis:2020unq,Ethier:2021bye,Celada:2024mcf}, is therefore a statement about the
% truncation of the expansion, and it matters most where the effects of the operators are largest, at
% high energies. The same growth with the energy that makes the tails of kinematic distributions
% sensitive to new physics thus also sets the limit of validity of the
% expansion~\cite{Contino:2016jqw}, a point to which we return below.

\subsubsection*{Operator Bases}

Writing down the operators is less straightforward than it may appear. Different operators can give
identical contributions to every observable, since they are related by integration by parts, by field
redefinitions~\cite{Politzer:1980me,Georgi:1991ch,Arzt:1993gz,Criado:2018sdb}, by Bianchi and
group-theoretic identities, and, for fermions, by Fierz identities. Moreover, in four dimensions, any
tensor that is anti-symmetric in five Lorentz indices vanishes. This yields identities involving the
Levi-Civita tensor, which relate operators with few fields, such as those containing dual field
strengths, and implies that Gram determinants built out of five independent momenta vanish, which,
because of momentum conservation, constrains operators with at least six
fields~\cite{Henning:2017fpj}. 

A basis is a maximal set of operators that remain independent once
all these relations are imposed~\cite{Lehman:2015via,Henning:2015alf}, and only in such a basis are
the Wilson coefficients uniquely defined. An equivalent characterization, used throughout this
thesis, identifies each independent operator with an independent contact amplitude, i.e., with a
local kinematic structure consistent with the helicities and the internal quantum numbers of the
external
particles~\cite{Shadmi:2018xan,Ma:2019gtx,Durieux:2019eor,Durieux:2020gip,Li:2020gnx,Dong:2021vxo,Li:2022tec}.
This amplitude-operator correspondence makes the redundancies of the Lagrangian transparent. Indeed,
operators that differ by terms vanishing on shell or by a field redefinition generate the same
on-shell amplitudes, while integration by parts translates into momentum conservation. Furthermore,
the relations specific to four dimensions, such as Fierz identities and those implied by vanishing
Gram determinants, are built into the spinor-helicity variables, where they reduce to Schouten
identities among spinor brackets.

The
SMEFT~\cite{Buchmuller:1985jz,Grzadkowski:2010es,Brivio:2017vri,Isidori:2023pyp,Aebischer:2025qhh,Falkowski:2023hsg,Henning:2014wua}
is the realization of Eq.~\eqref{eq:intro_Leff} built out of the SM fields and invariant
under its gauge group, with the Higgs field in an $SU(2)_L$ doublet, so that the electroweak
symmetry is linearly realized. Its only operator at dimension five, up to flavor indices, is the
Weinberg operator, which violates lepton number by two units and generates Majorana masses for the
neutrinos after electroweak symmetry breaking~\cite{Weinberg:1979sa}. At dimension six, the
operators conserving baryon number reduce to $59$ independent structures for one generation of
fermions, collected in the Warsaw basis~\cite{Grzadkowski:2010es}, which become $2499$ independent
real parameters once the flavor structure of three generations is taken into
account~\cite{Alonso:2013hga}. The SMEFT assumes that no light states other than those of the
SM exist below $\Lambda$, and it has become the reference framework for the combined
interpretation of collider, electroweak, and flavor
data~\cite{Brivio:2019ius,Ellis:2020unq,Ethier:2021bye,Celada:2024mcf}.

That assumption fails if new light states are present, and a well-motivated example is the QCD
axion. The QCD Lagrangian admits a CP-violating term proportional to $G\widetilde
G$~\cite{tHooft:1976rip,Callan:1976je,Jackiw:1976pf}, whose physical coefficient $\overline\theta$
combines the vacuum angle of QCD with the phase of the determinant of the quark mass matrix. This
term induces an electric dipole moment of the neutron~\cite{Baluni:1978rf,Crewther:1979pi}, whose
experimental bound~\cite{Abel:2020pzs} requires $|\overline\theta|\lesssim10^{-10}$, a value for
which the SM offers no explanation. This is the strong CP problem. Peccei and Quinn
proposed to solve it by introducing a global $U(1)$ symmetry, anomalous under QCD and spontaneously
broken at a scale $f_a$~\cite{Peccei:1977hh,Peccei:1977ur}. The corresponding
pseudo-Nambu--Goldstone boson, the axion~\cite{Weinberg:1977ma,Wilczek:1977pj}, couples to
$G\widetilde G$ with a strength proportional to $1/f_a$, so that $\overline\theta$ is effectively
promoted to a dynamical field. Moreover, the non-perturbative dynamics of QCD generates a potential
for the axion, whose minimum sets the effective $\overline\theta$ to zero~\cite{Vafa:1984xg}, and
gives it a mass of order $m_\pi
f_\pi/f_a$~\cite{Weinberg:1977ma,Wilczek:1977pj,GrillidiCortona:2015jxo}. The original proposal,
with $f_a$ at the electroweak scale, was soon excluded
experimentally~\cite{Donnelly:1978ty,DiLuzio:2020wdo}, but models with $f_a$ well above
it~\cite{Kim:1979if,Shifman:1979if,Dine:1981rt,Zhitnitsky:1980tq} remain viable. In this case, the
axion is long-lived and can be produced non-thermally in the early universe, which makes it a
natural dark-matter candidate~\cite{Abbott:1982af,Preskill:1982cy,Dine:1982ah,Davis:1986xc}.

ALPs generalize the axion to any light pseudoscalar $\phi$ that arises as the
pseudo-Nambu--Goldstone boson of an approximate global symmetry, spontaneously broken at a scale
well above the electroweak one~\cite{Jaeckel:2010ni,Marsh:2015xka,Irastorza:2018dyq,Choi:2020rgn}.
Their lightness is protected by the symmetry itself, but, unlike for the QCD axion, their mass and
the scale of symmetry breaking are independent parameters, since the explicit breaking of the
symmetry need not originate from QCD. ALPs bear on several of the open questions listed above,
including the flavor
puzzle~\cite{Davidson:1981zd,Wilczek:1982rv,Berezhiani:1989fp,Calibbi:2016hwq,Ema:2016ops,Greljo:2024evt},
the stability of the electroweak scale~\cite{Graham:2015cka}, and the nature of dark
matter~\cite{Arias:2012az}. The EFT of ALPs adds $\phi$ to the SM field content, and its
Goldstone origin shapes the way it couples to the SM. Indeed, a shift symmetry
$\phi\to\phi+c$ allows couplings only through the derivative $\partial_\mu\phi$ or through the
products of $\phi$ with $G\widetilde G$, $W\widetilde W$, and $B\widetilde B$, which change by a
total derivative, so that the symmetry is
broken only by non-perturbative effects, such as those that generate the potential of the QCD axion.
The leading interactions then arise at dimension five~\cite{Georgi:1986df} and are suppressed by the
symmetry-breaking scale. The operators invariant under the shift symmetry are known up to dimension
eight~\cite{Song:2023lxf,Grojean:2023tsd,Bertuzzo:2023slg}. If the symmetry is also broken
explicitly, couplings such as $\phi GG$, $\phi WW$, $\phi BB$, and non-derivative couplings to
fermions become possible, and, since they have CP properties opposite to those of the shift-symmetric
ones, their simultaneous presence violates
CP~\cite{Moody:1984ba,Marciano:2016yhf,DiLuzio:2020oah,DiLuzio:2023cuk,DasBakshi:2023lca}.

The phenomenology of ALPs raises both of the questions addressed in this thesis. On the one hand,
ALPs are searched for over a wide range of energies, from astrophysical and cosmological
observations and low-energy experiments to flavor factories and high-energy
colliders~\cite{Irastorza:2018dyq,Bauer:2017ris,Bauer:2021mvw}, so that couplings defined at the
scale of symmetry breaking must be evolved down to very different scales. Along this evolution, the
couplings mix, and, in particular, the couplings to gauge bosons induce couplings to
fermions, whose effects require a flip of the
fermion chirality and are therefore proportional to the fermion masses. On the other hand, all the
ALP interactions have dimension five or higher, so that the amplitudes they generate grow with the
energy, and searches at high-energy
colliders~\cite{Bauer:2017ris,Brivio:2017ije,Gavela:2019cmq} probe a regime in which the consistency
of the effective description must be verified~\cite{Brivio:2021fog}.

\subsubsection*{Matching and Running}

Whichever approach is adopted, the Wilson coefficients are tied to the UV in two steps, illustrated
in Figure~\ref{fig:intro_matching_running} for new physics above the electroweak scale. In the
first, known as matching, the heavy states are integrated out at a scale $\mu\sim\Lambda$, and the
coefficients are fixed by requiring the EFT to reproduce the low-energy amplitudes of the full
theory. In the top-down approach, this determines the $c_i(\Lambda)$ as functions of the couplings
and masses of a specific UV model. 

In the second step, known as running, the coefficients are evolved from $\mu\sim\Lambda$ down to the
scale $\mu\sim E$ of the measurement, according to the renormalization-group
equations~\cite{Callan:1970yg,Symanzik:1970rt}
\begin{equation}\label{eq:intro_RGE}
    \mu\frac{\dd c_i}{\dd\mu} = \sum_j \gamma_{ij}\,c_j + \dots\,,
\end{equation}
where $\gamma_{ij}$ is the anomalous-dimension matrix, which depends on the couplings and mass
parameters of the theory, and the ellipsis denotes terms of higher order in the Wilson coefficients,
arising from multiple insertions of higher-dimensional operators. Solving these equations resums the
large logarithms of $\Lambda/E$ that would otherwise spoil perturbation theory. To first order in
the logarithm, e.g.,
\begin{equation}\label{eq:intro_LL}
    c_i(E) \simeq c_i(\Lambda) - \sum_j\gamma_{ij}\,c_j(\Lambda)\,\log\frac{\Lambda}{E}\,,
\end{equation}
while the full solution resums all powers of $\gamma_{ij}\log(\Lambda/E)$. The diagonal entries of
$\gamma_{ij}$ describe the multiplicative renormalization of each operator, while the off-diagonal
ones describe operator mixing, by which a coefficient generated at $\Lambda$ induces, at lower
energies, coefficients that were absent at the matching scale. Mixing is not restricted to operators
of the same dimension. Indeed, since in mass-independent schemes the anomalous dimensions depend
polynomially on the masses, a higher-dimensional operator can mix into lower-dimensional ones
through positive powers of the mass parameters of the theory, while products of lower-dimensional
coefficients can source higher-dimensional
operators~\cite{Jenkins:2013zja,Chala:2021pll,DasBakshi:2022mwk,Bresciani:2023jsu}. The two steps
alternate whenever a threshold is crossed. In particular, at the electroweak scale, the $W$ and $Z$
bosons, the Higgs boson, and the top quark are integrated out, and the SMEFT is matched onto the
Low-Energy Effective Field Theory (LEFT)~\cite{Jenkins:2017jig,Dekens:2019ept}, whose coefficients
$\widehat c_a$ are then run down to the scale of the measurement. Running is
the central subject of the second part of this thesis.

\begin{figure}[htbp]
\centering
\begin{tikzpicture}[
    font=\small,
    theory/.style={draw=blobcolor!80!black, fill=blobcolor!15, rounded corners=2pt, align=center,
                   minimum width=3.4cm, minimum height=0.95cm},
    observ/.style={draw=ffcolor!80!black, fill=ffcolor!15, rounded corners=2pt, align=center,
                   minimum width=3.4cm, minimum height=0.95cm},
    flow/.style={-{Latex[length=2.2mm]}, thick},
    eqn/.style={anchor=west, align=left, fill=white, inner sep=2pt},
]
    \draw[flow] (0,-1.3) -- (0,10.3) node[above] {$\mu$};
    \foreach \y/\lab in {8/{\Lambda}, 4.4/{m_W}, 0.8/{E}} {
        \draw[thick] (-0.12,\y) -- (0.12,\y);
        \node[left] at (-0.12,\y) {$\lab$};
        \draw[densely dashed, gray!70] (0.12,\y) -- (11.8,\y);
    }
    \node[theory] (uv)  at (2.4,9.3)  {UV theory\\[1pt] $\mathcal L_{\rm UV}$};
    \node[theory] (eft) at (2.4,6.2)  {SMEFT\\[1pt] $\mathcal L_{\rm SMEFT}(c_i)$};
    \node[theory] (low) at (2.4,2.6)  {LEFT\\[1pt] $\mathcal L_{\rm LEFT}(\widehat c_a)$};
    \node[observ] (obs) at (2.4,-0.45) {Observables at $E$};
    \draw[flow] (uv.south)  -- (eft.north);
    \draw[flow] (eft.south) -- (low.north);
    \draw[flow] (low.south) -- (obs.north);
    \draw[flow, blobcolor!80!black, line width=1.2pt] (4.6,7.85) -- (4.6,4.55);
    \draw[flow, blobcolor!80!black, line width=1.2pt] (4.6,4.25) -- (4.6,0.95);
    \node[eqn] at (5.0,8) {\textbf{Matching}\quad
        $\mathcal M_{\rm UV}=\mathcal M_{\rm SMEFT}$ at $\mu\simeq\Lambda$
        $\;\Rightarrow\;c_i(\Lambda)$};
    \node[eqn] at (5.0,6.2) {\textbf{Running}\quad
        $\displaystyle\mu\frac{\dd c_i}{\dd\mu}=\sum_j\gamma_{ij}\,c_j+\dots$\\[6pt]
        $\displaystyle c_i(\mu)\simeq c_i(\Lambda)-\sum_j\gamma_{ij}\,c_j(\Lambda)\log\frac{\Lambda}{\mu}+\dots$};
    \node[eqn] at (5.0,4.4) {\textbf{Matching}\quad
        $\mathcal M_{\rm SMEFT}=\mathcal M_{\rm LEFT}$ at $\mu\simeq m_W$
        $\;\Rightarrow\;\widehat c_a(m_W)$};
    \node[eqn] at (5.0,2.6) {\textbf{Running}\quad
        $\displaystyle\mu\frac{\dd \widehat c_a}{\dd\mu}=\sum_b\widehat\gamma_{ab}\,\widehat c_b+\dots$\\[6pt]
        $\displaystyle \widehat c_a(\mu)\simeq \widehat c_a(m_W)-\sum_b\widehat \gamma_{ab}\,\widehat c_b(m_W)\log\frac{m_W}{\mu}+\dots$};
    \node[eqn] at (5.0,0.8) {\textbf{Observables}\quad computed with $\widehat c_a(E)$};
\end{tikzpicture}
\caption[Matching and running from the UV scale to low-energy observables]{Schematic
representation of the matching and running steps that connect a UV theory at a scale $\Lambda$ above
the electroweak scale to observables at the energy $E$, through the SMEFT and the LEFT. At each threshold the heavy states are integrated out and the Wilson
coefficients are fixed by equating the amplitudes of the two descriptions, while between thresholds
they are evolved by the renormalization-group equations, which resum the large logarithms of the
ratios of scales.}
\label{fig:intro_matching_running}
\end{figure}

\subsubsection*{Anomalous Dimensions and Operator Mixing}

The first requirement, control over the evolution of the coefficients, therefore rests on the
anomalous-dimension matrix, which is a necessary ingredient for interpreting experimental results in
an EFT. Its off-diagonal entries determine how a bound obtained on one operator translates into
constraints on the coefficients of others, and they can make an observable sensitive to new physics
that does not contribute to it at the matching scale. Electric dipole moments and flavor-violating
decays are typical examples, since they are measured at energies of a few GeV or below with a
precision that allows them to probe, through mixing, operators generated at the TeV scale or
above. The one-loop anomalous-dimension matrix of the
dimension-six SMEFT is known in
full~\cite{Jenkins:2013zja,Jenkins:2013wua,Alonso:2013hga,Alonso:2014zka}, as is that of the
LEFT~\cite{Jenkins:2017jig,Jenkins:2017dyc}, and the running of the ALP couplings has been studied
at one loop~\cite{Bauer:2020jbp,Chala:2020wvs,DasBakshi:2023lca}.

The ALP offers a concrete illustration of why mixing matters. Since the simultaneous presence of its
CP-even and CP-odd couplings violates CP, pairs of such couplings mix at one loop into CP-odd
operators built out of SM fields, such as the three-gluon operator~\cite{Weinberg:1989dx}, the electric dipole operators of quarks and leptons, and the
chromo-electric dipole operators of quarks~\cite{Marciano:2016yhf,DiLuzio:2020oah,DiLuzio:2023cuk}.
These operators are in turn constrained by measurements of electric dipole moments at very low
energies. Therefore, how strongly these measurements constrain the original ALP couplings depends on
the logarithms generated by this mixing between the scale of symmetry breaking and the ALP mass, on
how the resulting operators run further down to the hadronic scale, and on how they mix with one
another along the way.

The one-loop computations in the SMEFT showed that anomalous-dimension matrices are much sparser
than dimensional analysis and power counting would suggest, since many mixings allowed by every
symmetry of the theory vanish at one loop. In the SMEFT the vanishing entries follow a pattern
reminiscent of the holomorphy of supersymmetric theories~\cite{Alonso:2014rga}, and several
mechanisms have since been identified behind them, from supersymmetric embeddings~\cite{Elias-Miro:2014eia} to selection rules based on helicity~\cite{Cheung:2015aba}, on angular momentum~\cite{Jiang:2020rwz}, and,
beyond one loop, on color~\cite{Bern:2020ikv} and length of the operators~\cite{Bern:2019wie}. 
However, most of these zeros are not manifest in the diagrammatic approach. Some follow from the absence of one-loop diagrams connecting the two operators. The others emerge only at the very end of the calculation, once all the diagrams contributing to a given mixing have been summed and, when the divergences also require counterterms that vanish on shell, once the result has been reduced from a Green's basis~\cite{Gherardi:2020det} to the physical one through field redefinitions. In this approach, these zeros therefore appear as cancellations among contributions that are individually non-zero and gauge dependent.
Instead, in the on-shell language, most of them follow
from simple arguments, as discussed below, and can be predicted before any loop computation is
performed.

Whether comparable structures persist beyond one loop is a question of direct practical relevance.
The separation between a UV theory at or above the TeV scale and the observables that constrain it
most strongly, such as dipole moments and rare flavor transitions, spans several orders of
magnitude, and the accuracy of the corresponding measurements is such that two-loop running is
required in a growing number of
analyses~\cite{Allwicher:2023aql,Stefanek:2024kds,Haisch:2024wnw,Olgoso:2025jot,Mantani:2026fao,Born:2026tgm,Haisch:2026eyq}.

More fundamentally, whenever an entry of the anomalous-dimension matrix vanishes at one loop, the
two-loop term is not a correction but the leading contribution to that entry,
and it competes with the iteration of one-loop mixings through intermediate operators. 
% The radiative decay $b\to s\gamma$ is a well-known instance. QCD corrections enhance its rate by a factor of about
% two to three, mainly through the mixing of four-quark operators into the photon dipole operator,
% which vanishes at one loop and first arises at two
% loops~\cite{Grinstein:1987vj,Misiak:1992bc,Ciuchini:1993ks,Buchalla:1995vs}, where it depends on the
% scheme and combines with one-loop finite terms into a scheme-independent result. 
The radiative decay $b\to s\gamma$ is a well-known instance, since the mixing of four-quark operators into the photon dipole operator vanishes at one loop, arises first at two loops, and is responsible for the QCD enhancement of the rate by a factor of about two to three~\cite{Grinstein:1987vj,Misiak:1992bc,Ciuchini:1993ks,Buchalla:1995vs}.
Therefore, knowing
which two-loop entries vanish, and which are nevertheless fixed by one-loop information, determines
which effects can be neglected and which must be computed, at a time when two-loop calculations in
EFTs~\cite{Aebischer:2022anv,Jenkins:2023bls,Jenkins:2023rtg,Born:2024mgz,DiNoi:2024ajj,Fuentes-Martin:2024agf,Naterop:2024cfx,Ibarra:2024tpt,Guedes:2025sax,Duhr:2025zqw,Duhr:2025yor,Naterop:2025lzc,Naterop:2025cwg,DiNoi:2025tka,DiNoi:2025arz,Banik:2025wpi,Aebischer:2025hsx,Zhang:2025ywe,Aebischer:2026zxv,Born:2026xkr} are becoming necessary to match the precision of current and upcoming measurements.
% The question is also subtler than at one loop. Within mass-independent schemes, and for a fixed
% basis of physical operators, the one-loop anomalous dimensions do not depend on the scheme, whereas
% from two loops onwards they do~\cite{Buras:1989xd,Dugan:1990df,Herrlich:1994kh}, since a finite
% renormalization at one loop feeds into the two-loop terms. In particular, a mixing that first arises
% at two loops is in general scheme dependent already at leading order, as in the case of $b\to
% s\gamma$. Moreover, beyond one loop, the one-loop amplitudes that feed into the two-loop mixing
% contain finite terms that are not constrained by the one-loop selection rules, and the question
% becomes which zeros survive and in which sense they are independent of the
% scheme.
The question is also subtler than at one loop since, from two loops onward, the anomalous dimensions depend on the renormalization scheme, and one must establish which zeros survive and in which sense they are independent of the scheme.

% Moreover, anomalous dimensions depend on the field content and on the gauge group
% of the theory at hand. For the renormalizable couplings this was addressed long ago, up to two
% loops, by computing the renormalization-group equations of a general gauge theory with arbitrary
% scalars and fermions, from which those of any specific model follow by group-theoretic
% manipulations~\cite{Machacek:1983tz,Machacek:1983fi,Machacek:1984zw,Jack:1984vj}.

Moreover, anomalous dimensions depend on the field content and on the gauge group of the theory at
hand, so that they would, in principle, have to be recomputed for every model. For the renormalizable
couplings, this is avoided by the renormalization-group equations of a general gauge theory with
arbitrary scalars and fermions, known up to two loops, from which those of any specific model follow
by group-theoretic manipulations
alone~\cite{Machacek:1983tz,Machacek:1983fi,Machacek:1984zw,Jack:1984vj}.

Extending this program to higher-dimensional operators is valuable for the same reason, and for a
more specific one as well. The SMEFT and the ALP EFT are both special cases of a theory with scalars
and fermions in arbitrary representations of a gauge group, together with its gauge bosons,
and so are many other extensions of the SM by light scalars and fermions. These light
states can leave distinctive signatures through their mixing into the well-constrained operators of
the SMEFT, as shown for the ALP in Ref.~\cite{Galda:2021hbr}. Therefore, a general framework makes
such effects accessible for any of these extensions once the group and the representations are
specified, without a new loop computation.

\subsubsection*{Constraints from First Principles}

The second requirement concerns the values that the Wilson coefficients may take. Because the operators
in Eq.~\eqref{eq:intro_Leff} have dimension larger than four, the amplitudes they generate grow with
the energy, and at some point they become too large to be compatible with the conservation of
probability. Indeed, unitarity of the $S$-matrix, projected onto states of definite total angular
momentum, bounds each partial-wave amplitude by a number of order one. Since the partial-wave
amplitude generated by a dimension-six operator in a $2\to2$ process grows at most as $c_i
s/\Lambda^2$, the bound translates into an upper limit on $|c_i|$ that decreases with the energy,
and the energy at which a given coefficient saturates it marks the scale at which the effective
description ceases to be perturbatively reliable. 

Applied to the scattering of longitudinally
polarized $W$ bosons, the same reasoning led Lee, Quigg, and Thacker to bound the Higgs mass at
about $1$~TeV or, had there been no Higgs boson, to place the onset of new dynamics at that
scale~\cite{Lee:1977yc,Lee:1977eg}. This argument, later developed into the no-lose
theorem~\cite{Chanowitz:1985hj}, was of paramount importance in motivating the construction of the LHC.

In the EFT context, partial-wave unitarity bounds serve a related
purpose~\cite{Gounaris:1994cm,Corbett:2014ora,Corbett:2017qgl,DiLuzio:2016sur,DiLuzio:2017chi,Falkowski:2019tft,Chang:2019vez,Cohen:2021ucp,Brivio:2021fog,Cao:2023qks}.
Searches for new physics at colliders increasingly rely on the high-energy tails of
kinematic distributions, and the same enhancement that makes those tails sensitive also makes the
truncation of Eq.~\eqref{eq:intro_Leff} questionable there~\cite{Contino:2016jqw}. Therefore, an interpretation of such
data in terms of Wilson coefficients requires a quantitative statement of where the effective
description remains consistent, and unitarity bounds supply it from first principles, independently
of any experimental input. Being tree-level statements, they do not exclude regions of parameter
space in the strict sense. Instead, they indicate the energy above which the EFT must be completed,
by new states or by strong dynamics, which is the information needed to assess the validity of an
EFT analysis.

The standard implementation of these bounds decomposes $2\to2$ helicity amplitudes into partial
waves by means of the Wigner $d$-matrices~\cite{Jacob:1959at}, and diagonalizes the matrix of
partial-wave amplitudes connecting all states with the same quantum numbers, imposing the unitarity
condition on its eigenvalues. This coupled-channel analysis is necessary, because coefficients that
enter the same partial wave are correlated and must be constrained together. Since the partial-wave
coefficients depend linearly on the Wilson coefficients, the unitarity conditions carve out a convex
region in coefficient space, and the bound on a single coefficient obtained by letting all the
others vary, known as a marginalized bound, is a projection of that region, conservative by
construction and distinct from the bound obtained by switching on one operator at a time. However,
the method is tied to the kinematics of two-body states, and earlier treatments of processes with
more particles were restricted to the $J=0$ projection of contact
amplitudes~\cite{Chang:2019vez,Falkowski:2019tft,Abu-Ajamieh:2020yqi,Cohen:2021ucp}. Processes with
more particles in the final state, whose partial waves can grow faster with energy than those of
the $2\to2$ processes induced by the same operator, for instance when the latter require an
insertion of the Higgs vacuum expectation value, are expected to dominate at the energies of the
High-Luminosity LHC and of future colliders~\cite{FCC:2018byv,Accettura:2023ked}, and therefore lie
largely outside it. This matters most for operators that generate no energy-growing $2\to2$
amplitude at all, the six-Higgs operator $(H^\dagger H)^3$ of the SMEFT being the
simplest example. Moreover, the same questions arise
for EFTs of gravity, whose higher-derivative interactions parametrize the corrections to general
relativity at low energies~\cite{Endlich:2017tqa,Ruhdorfer:2019qmk}, and for which, even for $2\to2$
processes, the construction becomes impractical because amplitudes of particles with spin two are
cumbersome to evaluate from Feynman rules.

Unitarity is not the only property of the $S$-matrix that constrains the Wilson coefficients.
Indeed, analyticity, which follows from causality, together with crossing symmetry, relates the
low-energy amplitude to its behavior at all energies through dispersion relations. For forward
elastic scattering, the coefficient of $s^2$ at low energies is thereby expressed as an integral
over the total cross sections of the process and of its crossed counterpart, which are positive,
provided the amplitude grows more slowly than $s^2$ at high energies~\cite{Froissart:1961ux}. 

When
the EFT is the low-energy limit of a unitary, local, and causal completion, this yields inequalities
on specific combinations of coefficients, known as positivity
bounds~\cite{Pham:1985cr,Ananthanarayan:1994hf,Adams:2006sv}. They are derived in
Section~\ref{sec:analyticity} and apply most directly to operators whose forward amplitudes grow as
$s^2$, which in the theories of interest arise at dimension eight, the quartic photon interactions
of the Euler--Heisenberg Lagrangian~\cite{Heisenberg:1936nmg} being a familiar example. 

At dimension
six, the forward dispersion relation does not fix a sign. However, for four-fermion operators,
spin-dependent dispersion relations, known as spinning sum rules, relate the sign of certain
combinations of coefficients to the spin of the heavy states exchanged in the
UV~\cite{Remmen:2020uze,Remmen:2022orj,Gu:2020thj,Azatov:2021ygj}, at the price of additional
assumptions on the UV completion. Positivity and unitarity bounds are complementary in character,
since the former confine combinations of coefficients to a cone, restricting their signs, and the
latter bound their magnitudes, so that the region allowed by both can be considerably smaller than
either alone.

% These three classes of constraints rest on different assumptions, which is worth keeping in mind
% whenever they are compared or combined. Partial-wave unitarity bounds require only that the EFT
% remain perturbative at the energy considered. Positivity
% bounds do not require perturbativity, but require that some unitary, local, and causal completion
% exists, without specifying it. Spinning sum rules further assume a tree-level completion dominated
% by states of definite spin, together with a high-energy behavior of the amplitude soft enough for
% the contributions from infinity to vanish, which fails for instance when heavy vectors are exchanged
% in the $t$ channel~\cite{Azatov:2021ygj}. In exchange for these assumptions they return a different
% kind of information, namely whether the new dynamics is dominated by scalar or by vector states.
% Therefore, a coefficient lying outside the regions selected by the sum rules does not exclude a UV
% completion, but points to a more structured one, involving for instance vectors exchanged in the $t$
% channel, loop effects, or a mixture of spins.



The practical interest of such constraints grows with the size of the parameter spaces involved.
With about $2500$ independent coefficients at dimension six, the SMEFT has far more parameters than
current data can constrain simultaneously, and global
analyses~\cite{Brivio:2019ius,Ellis:2020unq,Ethier:2021bye,Celada:2024mcf} benefit from criteria
that restrict the viable region without assuming a specific model. Constraints derived from first
principles serve this purpose, and can be incorporated into such analyses, for instance as priors
that restrict the parameter space to physically allowed regions, and into the generation of
simulated events.

\subsubsection*{On-Shell Methods}

The two programs outlined above share a common technical obstacle. Computations in EFTs become
quickly intricate, since the number of Feynman diagrams grows rapidly with the number of loops and
external legs, and higher operator dimensions, multiple operator insertions, and external particles
with spin each make the problem worse. The difficulty is not only one of size. Indeed, individual
diagrams depend on the choice of gauge and on how the fields are continued off shell, only their sum
being physical, and in the computation of anomalous dimensions the divergences must first be
absorbed in a Green's basis and then projected onto the physical one through field redefinitions, as
noted above.

On-shell methods remove these intermediate steps by never leaving the space of physical amplitudes.
Their ingredients are gauge-invariant functions of the external momenta and helicities, constrained
by unitarity, locality, and
analyticity~\cite{Mangano:1990by,Dixon:1996wi,Elvang:2013cua,Dixon:2013uaa,Cheung:2017pzi,Travaglini:2022uwo}.

They are particularly efficient at high energies. When the energy of a process is much larger than
the masses of the particles involved, the latter can be neglected, and the amplitudes are those of
massless particles, labeled by their momenta, their internal quantum numbers, and their helicities.
In this regime, the spinor-helicity
variables~\cite{Berends:1981rb,DeCausmaecker:1981jtq,Gunion:1985vca,Xu:1986xb} make the amplitudes
particularly compact~\cite{Parke:1986gb}, and Lorentz invariance, locality, and unitarity are
restrictive enough to fix them to a large extent. For instance, for complex momenta, the three-particle
amplitudes are determined by the helicities of the particles up to a coupling
constant~\cite{Benincasa:2007xk}, and amplitudes with more particles follow from their factorization
on poles and cuts~\cite{Britto:2004ap,Britto:2005fq}, up to contact terms, which in an EFT
correspond to the independent operators. 
Moreover, helicity is a Lorentz-invariant label of massless
states, so that it can be tracked through every step of a calculation and gives rise to selection
rules. This is the regime relevant for this thesis. 

On the one hand, UV divergences are local, and
in mass-independent schemes such as $\MSbar$~\cite{tHooft:1973mfk,Weinberg:1973xwm,Bardeen:1978yd},
the anomalous dimensions therefore depend on the masses only polynomially, as on any other coupling,
so that they can be computed from massless amplitudes with the masses treated as perturbative
insertions. On the other hand, unitarity bounds become relevant at energies where the amplitudes
generated by higher-dimensional operators grow large, far above the masses of the SM
particles. Finally, massive particles can be treated within the same framework through the massive
spinor-helicity formalism~\cite{Arkani-Hamed:2017jhn}, 
% which makes the little group of massive
% particles manifest and has proven powerful as well, for instance in the construction of the
% electroweak EFT directly from on-shell amplitudes~\cite{Durieux:2019eor}.
which makes their little group manifest and extends the advantages of the massless formalism to massive states.

Unitarity relates the discontinuity of an amplitude to products of on-shell amplitudes of lower loop
order~\cite{Bern:1994zx,Bern:1994cg,Britto:2004nc}. Indeed, a unitarity cut puts the intermediate
states on shell, so that the discontinuity is assembled from lower-order physical amplitudes alone.
The same holds for the form factors of the operators, i.e., for their matrix elements between the
vacuum and on-shell multi-particle states. Here, the discontinuity has a direct physical meaning,
since the action of the dilatation operator on a form factor is fixed by the phase of the
$S$-matrix, so that anomalous dimensions can be read off from unitarity
cuts~\cite{Caron-Huot:2016cwu,EliasMiro:2020tdv,Baratella:2020lzz,Jiang:2020mhe,Bern:2020ikv}. The
one-loop renormalization-group equations of the second part of this thesis are obtained in this way,
as two-particle phase-space integrals of products of tree-level amplitudes and form factors, while,
at two loops, one-loop amplitudes and form factors enter the cuts as well. 

Read differently, the same
unitarity relation gives the bounds of the third part. Projected onto states of definite angular
momentum, it bounds the partial-wave coefficients, while its discontinuity, combined with
analyticity through a dispersion relation, confines the Wilson coefficients to a convex cone.

The formalism also has features that shape the results of the second part, one practical and one
limiting. Unitarity-based methods are most naturally formulated for massless particles, whereas many
phenomenologically relevant mixings, for instance, between chirality-preserving and
chirality-violating operators below the electroweak scale, require insertions of fermion masses. The
Higgs low-energy theorem~\cite{Ellis:1975ap,Shifman:1979eb} allows these leading mass effects to be
captured by insertions of a zero-momentum Higgs boson~\cite{Balkin:2021dko,Bertuzzo:2023slg}, so
that the formalism can remain massless throughout. 
On the other hand, four-dimensional unitarity
cuts do not determine the finite rational terms of one-loop amplitudes and form factors, which have
no discontinuity~\cite{Bern:1994zx,Bern:1994cg}. This is what divides the two-loop analysis into
zeros that are fixed by tree-level data and hold in any scheme that preserves the selection rule
based on the length of the operators, discussed below, and zeros that hold only in a suitably chosen
scheme.

Within these limits, working with physical amplitudes makes manifest structures that are hidden in the Lagrangian formulation. Following
Ref.~\cite{Cheung:2015aba}, one can assign to each operator a pair of weights $(w,\overline
w)=(\ell-h,\ell+h)$, where the length $\ell$ is the number of particles in the minimal amplitude
the operator generates and $h$ is their total helicity. The weights of the amplitudes entering a
unitarity cut combine in a definite way, and the tree-level amplitudes of the renormalizable theory
with four or more legs have $w,\overline w\ge4$, 
% unless a complex scalar couples to fermion pairs of
% the same chirality both directly and through its conjugate, as the Higgs doublet does in the
% SM. 
unless non-holomorphic Yukawa interactions are present.
Therefore, at one loop, an operator renormalizes only operators whose weights are at
least as large as its own, up to a shift of two units in the presence of such non-holomorphic Yukawa
couplings. Moreover, at one loop, no operator renormalizes a shorter one~\cite{Bern:2019wie}, so
that, for instance, four-fermion operators cannot renormalize the operators built from three field
strengths, and much of the sparsity of anomalous-dimension matrices discussed above becomes a
consequence of counting weights and lengths. 

Furthermore, operators that differ only in their CP
properties, such as $\phi GG$ and $\phi G\widetilde G$, generate the same on-shell helicity amplitudes up to a phase, so that their anomalous dimensions are manifestly related. 
This
relation is hidden in the diagrammatic approach, where such operators have Feynman rules with completely different
Lorentz structures. 

For the unitarity bounds, the on-shell approach extends the analysis to any process. Since the angular-momentum
content of an amplitude can be read off from its spinor-helicity representation, the partial-wave
decomposition can be constructed directly, for any number of external particles and any spin, by
diagonalizing a Casimir operator of the Poincaré group, the square of the Pauli--Lubanski vector, on
independent kinematic structures~\cite{Jiang:2020rwz,Shu:2021qlr}. Together with the
amplitude-operator correspondence, this ensures that, once processes of all the relevant
multiplicities are included, the unitarity conditions leave no flat direction in the space of Wilson
coefficients, so that the marginalized bounds remain finite.

The on-shell formalism required in the following chapters is collected in
Chapter~\ref{ch:amplitudes}. It introduces the spinor-helicity variables and the polarizations of
massless states up to spin two, discusses unitarity, locality, and analyticity of the $S$-matrix,
including the derivation of the dispersion relations and positivity bounds used in the third part,
and illustrates how these principles fix three-particle amplitudes and determine higher-point ones
recursively. 

% The remaining technical ingredients, such as form factors, the relation between the
% dilatation operator and the $S$-matrix, and the angular-momentum basis, are introduced in the
% chapters where they are first needed.

\subsubsection*{Outline of the Thesis}

The thesis is organized in four parts. The first, of which this Introduction is the opening chapter,
sets out the motivation and the tools, and continues with the review of on-shell methods in
Chapter~\ref{ch:amplitudes}. The second and third parts contain the original results. Each of their
chapters is based on one or two articles, as indicated at its opening, is written so as to be
readable on its own, and, where needed, collects its technical material in appendices at its end.
The fourth part summarizes the results and discusses possible extensions.

The second part is devoted to the anomalous dimensions of EFTs, and consists of the following
chapters:
\begin{itemize}
    \item Chapter~\ref{ch:RGEmethod}, based on Ref.~\cite{Bresciani:2023jsu}, reviews the relation between the dilatation operator and the phase of the $S$-matrix and derives from it master
    formulae for the anomalous dimensions at one and two loops, valid for the most general mixing
    among operators of different dimensions and for multiple operator insertions. It also shows how
    leading mass effects can be included by exploiting the Higgs low-energy theorem and keeping the formalism massless. These formulae are the common tool of the three chapters that follow.
    \item Chapter~\ref{ch:ALP_RGEs}, based on Ref.~\cite{Bresciani:2024shu}, applies them to the
    one-loop renormalization of a CP-violating ALP coupled to SM fields, above and below
    the electroweak scale, including the dimension-five and -six operators generated once the ALP is
    integrated out, and compares the on-shell computation with a parallel one based on Feynman
    diagrams.
    \item Chapter~\ref{ch:generalEFT}, based on Refs.~\cite{Aebischer:2025zxg,Aebischer:2025ddl},
    turns to a general effective gauge theory. We construct its physical operator basis up to
    dimension six and derive the complete one-loop renormalization-group equations for all its
    couplings at order $1/\Lambda^2$, from which those of a specific theory follow by group-theoretic
    manipulations alone. The SMEFT, LEFT, and $O(n)$ scalar model serve both as cross-checks and as worked examples.
    \item Chapter~\ref{ch:2looptheorems}, based on Ref.~\cite{Bresciani:2026yuw}, asks which
    structure of the anomalous-dimension matrix survives beyond one loop. We show that helicity
    selection rules fix to zero, from tree-level data alone, the $L$-loop mixings of operators into
    operators $L-1$ units shorter whenever their weights differ sufficiently, a result that holds in
    every scheme. For targets at least as long as the source,
    we construct a scheme in which further two-loop zeros hold, and show that in any other scheme the
    same entries are determined by one-loop data. Finally, we apply these results to general operators of
    dimension five through eight.
\end{itemize}

The third part is devoted to theoretical constraints on the Wilson coefficients, and consists of the
following chapters:
\begin{itemize}
    \item Chapter~\ref{ch:generalizedUB}, based on Ref.~\cite{Bresciani:2025toe}, reformulates the
    partial-wave decomposition in a vectorial formalism, so that the partial-wave coefficients become inner products between
    amplitudes and angular-momentum basis elements, the latter obtained by diagonalizing the square
    of the Pauli--Lubanski operator on spinor-helicity monomials. This extends partial-wave unitarity
    bounds to processes with arbitrary multiplicity and spin. We apply it to the quartic interactions
    of photons and gravitons, comparing the resulting bounds with positivity constraints, and to $2\to3$
    processes in the SMEFT, identifying the energies above which they supersede the $2\to2$ ones.
    \item Chapter~\ref{ch:ALP_bounds}, based on Ref.~\cite{Bresciani:2025ojh}, applies the formalism
    to the ALP EFT up to dimension eight, including couplings to leptons that explicitly break the $SU(2)_L$ symmetry, whose effects in charged-meson and $W$-boson decays grow with the energy, and compares the resulting unitarity bounds with collider searches and rare decays. It also derives the complete set of positivity bounds on the dimension-eight operators and discusses their interplay with the unitarity bounds.
    \item Chapter~\ref{ch:SMEFT_bounds}, based on Ref.~\cite{Bresciani:2026acy}, performs a coupled-channel analysis of the dimension-six SMEFT over all helicity, gauge,
    and flavor configurations, marginalizing each bound over the remaining coefficients. It derives
    spinning sum rules for four-fermion operators in the Warsaw basis together with a criterion to
    impose them in flavor space, extends the bounds to the
    minimal-flavor-violation and
    $U(2)^5$ hypotheses on the flavor structure, and compares them with
    current and projected experimental limits.
\end{itemize}

% A few remarks may help the reader navigate the material. Chapter~\ref{ch:RGEmethod} is a
% prerequisite for the rest of the second part, and Chapter~\ref{ch:2looptheorems} relies in
% particular on its two-loop master formula, while Chapter~\ref{ch:generalizedUB} is a prerequisite
% for the rest of the third part, whose positivity bounds also draw on Section~\ref{sec:analyticity}.
% Several chapters contain extensive collections of results intended for consultation more than for
% continuous reading. This is the case for the class-by-class derivation of the anomalous dimensions
% of the ALP EFT in Section~\ref{sec:ALPuvanomalousdim}, for the renormalization-group equations of
% the general effective gauge theory in Section~\ref{sec:gEFTresults}, for the tables of two-loop
% zeros in Section~\ref{sec:2loopapplications}, and for the bounds on flavorful four-fermion operators
% in Appendix~\ref{app:SMEFT_flavorresults}. 